\subsubsection{Problems}\label{subsec:sobolev-problems}

\begin{enumerate}[label=\thesection.\arabic*.,ref=\thesection.\arabic*,start=1,series=book-problems]

\item\label{ex:sobolev-difference-quotient-l-two-convergence} 设 $u\in H^1(\mathbb{R})$. 证明 $u'(x)=\lim_{h\to0}(u(x+h)-u(x))/h$ 在 $L^2$ 收敛的意义下成立.

\item\label{ex:sobolev-line-integration-by-parts} 设 $u,v\in H^1(\mathbb{R})$. 证明\[\int_{-\infty}^{\infty}u(x)v'(x)\,\mathrm{d}x=-\int_{-\infty}^{\infty}u'(x)v(x)\,\mathrm{d}x.\]

\item\label{ex:compact-support-zero-boundary-sobolev} 假设 $u\in H^k(\Omega)$ 在 $\Omega$ 中具有紧支集. 证明 $u\in H_0^k(\Omega)$.

\item\label{ex:sobolev-embedding-critical-counterexample} 证明 $H^1(\mathbb{R}^2)$ 不是 $C(\mathbb{R}^2)$ 的子集. Hint: 考察 $|\ln|\bm{x}||$ 的幂.

\item\label{ex:one-dimensional-h-one-product} 假设 $u,v\in H^1(\mathbb{R})$. 证明乘积 $uv$ 也属于 $H^1(\mathbb{R})$.

\item\label{ex:sobolev-embedding-cusp-failure} 证明, 若 $\partial\Omega$ 具有尖点, 则 Sobolev 嵌入定理一般不成立. Hint: 考察在尖点处具有奇性的函数.

\item\label{ex:sobolev-continuous-functions-vanish-at-infinity} 设 $u\in H^s(\mathbb{R}^m)$, $s>m/2$. 证明 $\lim_{\bm{x}\to\infty}u(\bm{x})=0$.

\item\label{ex:weak-sobolev-convergence-lower-order-strong} 设 $\Omega$ 满足 Theorem~\ref{thm:compact-successive-sobolev-orders} 的假设. 证明 $H^{k+1}(\Omega)$ 中每个弱收敛序列都在 $H^k(\Omega)$ 中强收敛.

\item\label{ex:ehrling-sobolev-equivalent-norms} 证明 Corollary~\ref{cor:ehrling-sobolev-equivalent-norms}.

\item\label{ex:zero-boundary-compact-sobolev-embedding-no-regularity} 设 $\Omega$ 有界 (不对 $\partial\Omega$ 作假设). 证明 $H_0^{k+1}(\Omega)$ 紧嵌入 $H^k(\Omega)$.

\item\label{ex:continuous-function-trace-support-lemma-proof} 证明 Lemma~\ref{lem:continuous-function-trace-and-support}.

\item\label{ex:trace-without-half-order-loss-impossible} 证明不存在将 $H^k(\mathbb{R}^m)$ 映到 $H^k(\mathbb{R}^{m-1})$ 的连续迹算子.

\item\label{ex:one-direction-bounded-domain-weak-dirichlet-existence} 假设 $\Omega$ 在一个方向上有界. 建立 Dirichlet 问题 $\Delta u=f$, $u|_{\partial\Omega}=0$ 的弱解的存在性结果.

\item\label{ex:distribution-local-negative-sobolev-order} 设 $D$ 为有界开集, 满足 $\overline D\subset\Omega$, $f\in\mathcal{D}^{\prime}(\Omega)$. 证明存在整数 $k$, 使 $f|_D\in H^{-k}(D)$. Hint: 使用 Lemma~\ref{def:distribution} 及 Sobolev 嵌入定理.

\item\label{ex:mean-zero-sobolev-poincare-inequality} 假设 $\Omega$ 有界, 连通且具有 $1$ 阶延拓性质. 令\[V=\left\{u\in H^1(\Omega)\ \middle|\ \int_{\Omega}u(\bm{x})\,\mathrm{d}\bm{x}=0\right\}.\]证明存在常数 $C$, 使\[\int_{\Omega}|u(\bm{x})|^2\,\mathrm{d}\bm{x}\leq C\int_{\Omega}|\nabla u(\bm{x})|^2\,\mathrm{d}\bm{x}\quad\text{for all }u\in V.\]对 $H^1(\Omega)$ 的哪些其他子集 $V$, 也有这样的不等式成立?

\item\label{ex:l-two-into-negative-one-compact-embedding} 证明, 若 $\Omega$ 有界, 则 $L^2(\Omega)$ 紧嵌入 $H^{-1}(\Omega)$.

\item\label{ex:whole-space-poincare-no-uniform-constant} 证明不存在常数 $C$, 使\[\int_{\Omega}|u(\bm{x})|^2\,\mathrm{d}\bm{x}\leq C\int_{\Omega}|\nabla u(\bm{x})|^2\,\mathrm{d}\bm{x}\quad\text{for all }u\in H^1(\mathbb{R}^m).\]

\item\label{ex:narrow-unbounded-domain-smooth-not-l-p-example} 令\[\Omega=\{(x_1,x_2)\in\mathbb{R}^2\mid0<x_1<\infty,\quad0<x_2<e^{-x_1}\}.\]证明, 对每个 $p\in(2,\infty)$, 存在函数 $u_p$, 使对所有 $k\geq0$ 都有 $u_p\in H^k(\Omega)$, 但 $u_p\notin L^p(\Omega)$.

\item\label{ex:titchmarsh-fourier-l-p-sobolev-embedding} Titchmarsh 的一个经典定理断言, 若 $p\in[1,2)$, 则 Fourier 变换将 $L^p(\mathbb{R}^m)$ 映入 $L^q(\mathbb{R}^m)$, 其中 $\dfrac{1}{p}+\dfrac{1}{q}=1$. 利用这一结果证明, $H^1(\mathbb{R}^3)$ 对所有 $p\in[2,6)$ 都连续嵌入 $L^p(\mathbb{R}^3)$. (Note: $H^1(\mathbb{R}^3)$ 也连续嵌入 $L^6(\mathbb{R}^3)$.)

\item\label{ex:periodic-sobolev-spaces-fourier-characterization} 定义 $\mathbb{R}$ 上周期函数的 Sobolev 空间, 并用 Fourier 级数刻画它们. 周期函数的 Sobolev 空间与 $[0,2\pi]$ 上的 Sobolev 空间有什么关系? Hint: 回顾 Problem~\ref{ex:periodic-fourier-uniform-convergence}.

\item\label{ex:boundary-smooth-functions-not-dense-open-set-example} 举出一个开集的例子, 使 $H^1(\Omega)\cap C^{\infty}(\overline\Omega)$ 在 $H^1(\Omega)$ 中不稠密.

\item\label{ex:k-diffeomorphism-definition-redundancies} 讨论 $k$ 阶微分同胚的定义中可能多余的条件.

\item\label{ex:manifold-sobolev-inner-products-equivalence} 验证由 \eqref{eq:manifold-sobolev-local-inner-product} 定义的所有内积都等价.

\item\label{ex:vandermonde-matrix-nonsingularity} 设 $A_{ij}=a_j^i$, $i,j=0,\ldots,K$, 其中 $a_j$ 为互不相同的实数. (采用约定 $0^0=1$.) 证明 $\det\bm{A}\ne0$.

\end{enumerate}
